\documentclass[10pt, spanish]{beamer}
\usepackage[utf8]{inputenc}
\usepackage[spanish]{babel}
\usepackage{amsmath, amssymb, amsthm}
\usepackage{graphicx}
\usepackage{hyperref}
\usepackage{booktabs}

% Tema y color
\usetheme{Madrid}
\usecolortheme{default}

% Datos de la presentación
\title[Cálculo I - Semana 4]{Cálculo I \\ Funciones exponencial, logarítmica, trigonométricas e inversas}
\author[Matihus A. Molina Larios]{Matihus Alberth Molina Larios}
\institute[UNMSM]{
	Universidad Nacional Mayor de San Marcos \\
	Facultad de Ciencias Matemáticas \\
	\texttt{matihus.molina@unmsm.edu.pe}
}
\date{\today}

% Redefinir entornos
\newtheorem{definicion}{Definición}
\newtheorem{teorema}{Teorema}
\newtheorem{corolario}{Corolario}
\newtheorem{ejemplo}{Ejemplo}
\newtheorem{nota}{Nota}

\setbeamertemplate{theorem}[ams style]

\begin{document}
	
	\begin{frame}
		\titlepage
	\end{frame}
	
	\begin{frame}{Contenido}
		\tableofcontents
	\end{frame}
	
	\section{Función exponencial}
	\subsection{Definición y propiedades}
	\begin{frame}{Función exponencial}
		\begin{definicion}
			Para $a>0$, $a\neq1$, la función exponencial de base $a$ es $f(x)=a^x$, definida para todo $x\in\mathbb{R}$.
		\end{definicion}
		\begin{teorema}[Propiedades]
			\begin{itemize}
				\item Dominio: $\mathbb{R}$.
				\item Rango: $(0,\infty)$.
				\item Si $a>1$, es estrictamente creciente; si $0<a<1$, estrictamente decreciente.
				\item $a^{x+y}=a^x a^y$, $a^{x-y}=a^x/a^y$, $(a^x)^y=a^{xy}$.
			\end{itemize}
		\end{teorema}
	\end{frame}
	
	\begin{frame}{Gráfica y aplicaciones}
		\begin{ejemplo}
			Graficar $f(x)=2^x$ y $g(x)=\left(\frac12\right)^x$.
		\end{ejemplo}
		\begin{proof}
			Ambas pasan por $(0,1)$. $2^x$ crece rápidamente, $\left(\frac12\right)^x$ decrece. El rango es $(0,\infty)$.
		\end{proof}
	\end{frame}
	
	\section{Función logarítmica}
	\subsection{Definición como inversa de la exponencial}
	\begin{frame}{Logaritmo}
		\begin{definicion}
			El logaritmo en base $a$ ($a>0$, $a\neq1$) es la función inversa de $a^x$:
			\[
			\log_a x = y \quad \Longleftrightarrow \quad a^y = x, \quad x>0.
			\]
		\end{definicion}
		\begin{teorema}[Propiedades]
			\begin{itemize}
				\item Dominio: $(0,\infty)$.
				\item Rango: $\mathbb{R}$.
				\item $\log_a(xy)=\log_a x+\log_a y$, $\log_a(x/y)=\log_a x-\log_a y$, $\log_a(x^r)=r\log_a x$.
				\item Cambio de base: $\log_a x = \frac{\log_b x}{\log_b a}$.
			\end{itemize}
		\end{teorema}
	\end{frame}
	
	\begin{frame}{Dominio y rango}
		\begin{ejemplo}
			Determinar dominio de $f(x)=\ln(x^2-4)$.
		\end{ejemplo}
		\begin{proof}
			Se requiere $x^2-4>0$, es decir $|x|>2$. Luego $\operatorname{Dom}(f)=(-\infty,-2)\cup(2,\infty)$. El rango es $\mathbb{R}$.
		\end{proof}
	\end{frame}
	
	\section{Funciones trigonométricas}
	\subsection{Definiciones y propiedades básicas}
	\begin{frame}{Funciones seno y coseno}
		\begin{definicion}
			Para un ángulo $x$ en radianes, $\sen x$ y $\cos x$ se definen mediante el círculo unitario.
		\end{definicion}
		\begin{teorema}[Propiedades]
			\begin{itemize}
				\item Dominio: $\mathbb{R}$.
				\item Rango: $[-1,1]$.
				\item Periodicidad: $\sen(x+2\pi)=\sen x$, $\cos(x+2\pi)=\cos x$.
				\item Identidad fundamental: $\sen^2 x+\cos^2 x=1$.
			\end{itemize}
		\end{teorema}
	\end{frame}
	
	\begin{frame}{Tangente y otras funciones}
		\begin{definicion}
			$\tan x = \frac{\sen x}{\cos x}$, $\cot x = \frac{\cos x}{\sen x}$, $\sec x = \frac{1}{\cos x}$, $\csc x = \frac{1}{\sen x}$.
		\end{definicion}
		\begin{ejemplo}
			Dominio de $\tan x$: $\mathbb{R}\setminus\{\frac{\pi}{2}+k\pi,\ k\in\mathbb{Z}\}$. Rango: $\mathbb{R}$.
		\end{ejemplo}
	\end{frame}
	
	\section{Funciones trigonométricas inversas}
	\subsection{Definiciones y dominios restringidos}
	\begin{frame}{Arcoseno, arcocoseno y arcotangente}
		\begin{definicion}
			Para obtener funciones inversas, restringimos los dominios:
			\begin{itemize}
				\item $\arcsen x$: dominio $[-1,1]$, rango $[-\pi/2,\pi/2]$.
				\item $\arccos x$: dominio $[-1,1]$, rango $[0,\pi]$.
				\item $\arctan x$: dominio $\mathbb{R}$, rango $(-\pi/2,\pi/2)$.
			\end{itemize}
		\end{definicion}
		\begin{teorema}
			$\arcsen(\sen x)=x$ para $x\in[-\pi/2,\pi/2]$, pero no para todo $x$.
		\end{teorema}
	\end{frame}
	
	\begin{frame}{Ejemplo: dominio y rango}
		\begin{ejemplo}
			Hallar dominio y rango de $f(x)=\arcsen(2x-1)$.
		\end{ejemplo}
		\begin{proof}
			Dominio: $-1\le 2x-1\le 1 \Rightarrow 0\le 2x\le 2 \Rightarrow 0\le x\le 1$. Rango: $[-\pi/2,\pi/2]$.
		\end{proof}
	\end{frame}
	
	\section{Bibliografía}
	\begin{frame}{Referencias}
		\begin{thebibliography}{9}
			\bibitem{spivak} Michael Spivak, \textit{Cálculo infinitesimal}, Editorial Reverté, 1998.
			\bibitem{stewart} James Stewart, \textit{Cálculo de una variable}, Cengage Learning, 2017.
			\bibitem{apostol} Tom M. Apostol, \textit{Cálculo}, Vol. I, Reverté, 2002.
			\bibitem{leithold} Louis Leithold, \textit{El Cálculo}, Oxford University Press, 1998.
		\end{thebibliography}
	\end{frame}
	
\end{document}